\section{Обоснование решения математической постановки задачи}

Полученная постановка представляет собой задачу параметрической
настройки распределённой системы, функционирующей в условиях случайных
сетевых задержек и отказов узлов, при вероятностных ограничениях на
показатели целостности \eqref{eq:integrity_final},
доступности \eqref{eq:availability_system}
и согласованности \eqref{eq:consistency_log}.

Поведение системы определяется случайными процессами
$\Delta_{ij}(t)$, $D_{ij}(t)$ и $S_{i,u}(t)$,
характеризующими задержки передачи сообщений,
потери пакетов и отказ узлов.
Поэтому показатели
$M\big(A_i^{\mathrm{r}}(t;\theta)\big)$,
$M\big(A_a^{\mathrm{w}}(t;\theta)\big)$
и
$M\big(L(t;\theta)\big)$
являются случайными величинами и зависят от параметров $\theta$
через вероятностные характеристики сети и инфраструктуры.

В общем случае аналитическое выражение зависимости функционала
$J(\theta)$ от параметров системы отсутствует.
Поэтому задача поиска оптимальной конфигурации
естественным образом относится к классу задач
стохастической оптимизации,
в которых качество конфигурации оценивается
по результатам наблюдений или имитационного моделирования~\cite{Shapiro2009}.

В практических задачах такого типа рассматривается конечное множество
кандидатных конфигураций
$
\Theta_{\mathrm{grid}} \subseteq \Theta_{\mathrm{adm}},
$
где $\Theta_{\mathrm{adm}}$ --- допустимая область параметров системы,
определяемая ограничениями корректности протокола
и архитектурными требованиями.
Допустимые диапазоны параметров
$(m,\tau_{\mathrm{fo}},\tau_{\mathrm{el}},\tau_{\mathrm{sp}},r)$
выводятся далее при обосновании выбора параметров системы.

Для фиксированной конфигурации параметров
$\theta \in \Theta_{\mathrm{grid}}$
значение функционала $J(\theta)$ является случайной величиной,
поскольку зависит от случайной реализации сетевых задержек
и отказов узлов.
Поэтому используется оценка методом выборочного среднего~\cite{Kim2011}:

\[
\widehat{J}_N(\theta)
=
\frac{1}{N}
\sum_{s=1}^{N}
J^{(s)}(\theta),
\]

где $J^{(s)}(\theta)$ --- значение функционала
на $s$-м независимом прогоне системы.

Итоговая конфигурация выбирается как

\[
\theta^{*}
=
\arg\min_{\theta \in \Theta_{\mathrm{grid}}}
\widehat{J}_N(\theta).
\]

Минимальность решения понимается следующим образом:
конфигурация $\theta^{*}$ имеет наименьшее значение
оценочного функционала $\widehat{J}_N$
среди всех рассмотренных допустимых кандидатов.

После определения общей методики поиска решения
необходимо обосновать выбор параметров,
входящих в вектор $\theta$,
и показать,
каким образом каждый из них влияет
на выполнение ограничений целостности,
доступности и согласованности системы,
а также на значение функционала $J(\theta)$.

\subsection{Обеспечение целостности данных}

Согласно определению журнала \eqref{eq:Li} (и его префикса $\mathcal{L}_i(t)$),
журнал кластера $c_i$ представляет собой упорядоченную последовательность записей
$\mathcal{L}_i=(e_{i,1},e_{i,2},\dots)$.

Межкластерная репликация определяется оператором доставленного префикса
$\mathcal{R}_{i\to j}(t;\theta)$,
который описывает часть журнала источника, принятую кластером $c_j$
к моменту времени $t$ при выбранной конфигурации параметров $\theta$.

В постановке задачи целостность \eqref{eq:integrity_final}
сводится к отсутствию события нарушения
$\mathsf{Int}_{i\to j}(t;\theta)$,
то есть к сохранению префиксности доставленного и применённого журнала.
Структурно это достигается тем, что записи применяются последовательно по номеру,
что предотвращает переупорядочивание и дублирование записей.

Однако указанное свойство не предотвращает возможность модификации
содержимого записи $e_{i,k}$ при передаче по каналу $(c_i,c_j)\in\mathcal{E}$.
Для обеспечения аутентичности и неизменности передаваемых данных
используется механизм криптографической аутентификации сообщений
на основе HMAC \cite{RFC2104}.

Для каждой записи $e_{i,k} \in \mathcal{L}_i$ формируется сообщение
\[
m_{i,k} = (i, k, e_{i,k}),
\]
где $i$ --- идентификатор исходного кластера,
$k$ --- логический номер записи (LSN),
$e_{i,k}$ --- содержимое записи журнала.
Для сообщения вычисляется код аутентичности
\[
\mathrm{tag}_{i,k}=\mathrm{HMAC}_{K}(m_{i,k}),
\]
где $K$ --- общий секретный ключ доверенных кластеров.

Передача по каналу осуществляется в виде пары $(m_{i,k},\mathrm{tag}_{i,k})$.
Получатель принимает запись к применению только при одновременном выполнении условий:
корректность $\mathrm{tag}_{i,k}$ и
соответствие номера $k$ ожидаемому значению, то есть $k = |\mathcal{L}_j(t)| + 1$.

С точки зрения постановки задачи это означает,
что при корректной работе механизма доставки и применения
выполняются отношения префикса
$\mathcal{R}_{i\to j}(t;\theta) \preceq \mathcal{L}_i(t)$ и
$\mathcal{R}_{i\to j}(t;\theta) \preceq \mathcal{L}_j(t)$,
то есть не возникает событие нарушения целостности $\mathsf{Int}_{i\to j}(t;\theta)$.

При корректной реализации HMAC и управлении ключами вероятность принятия
модифицированной записи является пренебрежимо малой,
что обеспечивает выполнение ограничения \eqref{eq:integrity_final}.

\subsection{Обоснование выбора параметра $m$}

Параметр $m$ входит в функционал $J(\theta)$ через вероятность сохранения
активности кластера $P_{\mathrm{act}}(m)$, введённую в разделе постановки.
В модели независимых отказов узлов (параметр среды $p$)
число исправных узлов имеет биномиальное распределение \cite{Kremer2008},
а вероятность того, что в активном кластере исправно не менее $m$ узлов, равна
\[
P_{\mathrm{act}}(m)
=
\sum_{k=m}^{n_a}
\binom{n_a}{k}
(1-p)^k p^{\,n_a-k}.
\]

В функционале \eqref{eq:J_final} учитывается величина $1-P_{\mathrm{act}}(m)$,
характеризующая риск потери активности кластера.
Увеличение параметра $m$ повышает требования к числу исправных узлов,
уменьшает $P_{\mathrm{act}}(m)$ и, следовательно, увеличивает слагаемое
$1-P_{\mathrm{act}}(m)$ в $J(\theta)$.

С точки зрения кворумных механизмов согласования
для корректного подтверждения операций требуется большинство узлов:
\[
m \ge \left\lfloor \frac{n_a}{2} \right\rfloor + 1.
\]
Это условие означает, что любые два множества узлов размера не менее $m$
имеют хотя бы один общий узел, что исключает принятие противоречивых решений
полностью разными группами узлов.

Таким образом, параметр $m$ задаёт компромисс между устойчивостью согласования
и вероятностью сохранения активности кластера и напрямую влияет на функционал $J(\theta)$,
а также косвенно влияет на выполнение ограничения доступности
\eqref{eq:availability_system} через величину $A_a^{\mathrm{w}}(t;\theta)$.

\subsection{Обоснование выбора параметров времени ожидания}

Параметры $\tau_{\mathrm{fo}}, \tau_{\mathrm{el}}, \tau_{\mathrm{sp}}$
влияют на скорость реакции системы на отказы и на вероятность ложных переключений.
В функционале \eqref{eq:J_final} они входят через отношения
\[
\frac{\tau_{\mathrm{el}}}{\tau_{\mathrm{fo}}}
\quad \text{и} \quad
\frac{\tau_{\mathrm{sp}}}{\tau_{\mathrm{fo}}},
\]
которые являются безразмерными и отражают относительные временные масштабы:
насколько быстро внутреннее восстановление и подтверждение состояния
происходят по сравнению с временем глобального переключения.

Рост $\tau_{\mathrm{el}}$ и $\tau_{\mathrm{sp}}$ увеличивает соответствующие слагаемые
в $J(\theta)$, поскольку внутренние процедуры занимают большую долю времени
по отношению к $\tau_{\mathrm{fo}}$.
С практической точки зрения это означает увеличение времени,
в течение которого система находится в промежуточных состояниях
и не может эффективно обслуживать запросы.

Для обеспечения корректной иерархии механизмов восстановления
используются соотношения, предотвращающие преждевременное межкластерное переключение:
\[
\tau_{\mathrm{fo}} \ge \alpha \tau_{\mathrm{el}}, \qquad \alpha>1,
\]
что гарантирует, что внутреннее восстановление (переизбрание лидера)
успевает завершиться до запуска глобального переключения,
и
\[
\tau_{\mathrm{sp}} \le \beta \tau_{\mathrm{fo}}, \qquad 0<\beta<1,
\]
что обеспечивает возможность выполнить несколько попыток подтверждения состояния
до принятия решения о смене роли кластера.

Данные условия уменьшают вероятность избыточных переключений и тем самым
способствуют выполнению ограничений доступности и согласованности,
то есть влияют на величины
$A_i^{\mathrm{r}}(t;\theta)$,
$A_a^{\mathrm{w}}(t;\theta)$
и
$L(t;\theta)$.

\subsection{Обоснование выбора параметра $r$}

Параметр $r$ входит в функционал $J(\theta)$ через вероятность успешной доставки
репликационного сообщения $P_{\mathrm{succ}}(r)$.
Если вероятность успешной доставки за одну попытку равна $p_d$,
а попытки передачи независимы, то
\[
P_{\mathrm{succ}}(r)=1-(1-p_d)^r.
\]
В функционале \eqref{eq:J_final} учитывается величина $1-P_{\mathrm{succ}}(r)$,
характеризующая риск полной потери сообщения после $r$ попыток передачи.
Увеличение $r$ монотонно увеличивает $P_{\mathrm{succ}}(r)$ и уменьшает
слагаемое $1-P_{\mathrm{succ}}(r)$, что способствует снижению лага репликации.

Одновременно рост $r$ увеличивает число сетевых передач и может повышать нагрузку,
что косвенно влияет на выполнение ограничений \eqref{eq:availability_system}
и \eqref{eq:consistency_log},
то есть на величины
$A_i^{\mathrm{r}}(t;\theta)$,
$A_a^{\mathrm{w}}(t;\theta)$
и
$L(t;\theta)$.
Следовательно, параметр $r$ определяет компромисс между снижением риска потери
сообщений и ростом эксплуатационной нагрузки системы, что отражается в значении $J(\theta)$.
